\section*{第 \ref{ch:b1:optical-instruments} 章——光学仪器与眼睛}

\begin{solution}{exo:b1:optical-instruments:1}
无穷远：$f' = \qty{17}{mm}$，$V = 1/0.017 = \qty{58.8}{\delta}$。
\qty{25}{cm}：$V = 58.8 + 1/0.25 = \qty{62.8}{\delta}$，$f' = \qty{15.9}{mm}$。
调节幅度 \qty{4.0}{\delta}。
\end{solution}

\begin{solution}{exo:b1:optical-instruments:2}
近视：无穷远处的物须成像于远点，$f' = \qty{-40}{cm}$，
$V = \qty{-2.5}{\delta}$。远视：位于 \qty{25}{cm} 处的物体须
成虚像于 \qty{1.0}{m} 处：$V = 1/\overline{OA'} - 1/\overline{OA}
= -1 + 4 = \qty{+3.0}{\delta}$。
\end{solution}

\begin{solution}{exo:b1:optical-instruments:3}
$G = 25/4.0 = 6.3$；邮票位于焦平面，
距透镜 \qty{4.0}{cm}。尺寸为 \qty{2}{mm} 的细节，其视角 $\theta' = 2/40 = \qty{0.05}{rad}$，
在 \qty{25}{cm} 处则为 $2/250 = \qty{0.008}{rad}$。
\end{solution}

\begin{solution}{exo:b1:optical-instruments:4}
$G = 900/25 = 36$ 和 $900/9 = 100$；筒长 $f_1' + f_2' =
\qty{925}{mm}$ 和 \qty{909}{mm}；月球视角：$36 \times 0.52^\circ = 19^\circ$ 和
$52^\circ$。
\end{solution}

\begin{solution}{exo:b1:optical-instruments:5}
$\gamma_1 = -160/4 = -40$；$G_2 = 250/20 = 12.5$；$G = 500$。
$\overline{F_1A} = -f_1'^2/\Delta = -16/160 = \qty{-0.10}{mm}$：物体
位于距物镜 \qty{4.1}{mm} 处。物体须与物镜保持该距离，
因此整个镜筒移动 \qty{10}{\micro m}——
这正是细调旋钮的作用。
\end{solution}

\begin{solution}{exo:b1:optical-instruments:6}
裸眼：$1.5/3\times10^{-4} = \qty{5}{km}$。望远镜：$\theta_{\min}
= 1.22 \times 5.5\times10^{-7}/0.100 = \qty{6.7e-6}{rad}$；两车灯
可分辨的最大间距为 $1.5/6.7\times10^{-6} \approx \qty{220}{km}$（原则上，
空气和曲率首先起作用）。$G_r = 3\times10^{-4}/
6.7\times10^{-6} = 45$。
\end{solution}

\begin{solution}{exo:b1:optical-instruments:7}
$D = 50/2 = \qty{25}{mm}$。曝光时间 $\propto N^2$，故为原来的 $(8/2)^2 = 16$ 倍，
即 $16/500 \approx 1/30$~\unit{s}。它换来了更大的景深（
并减少像差），代价是产生运动模糊，或需要三脚架。
\end{solution}

\begin{solution}{exo:b1:optical-instruments:8}
$H = 35^2/(11 \times 0.03) = \qty{3.7}{m}$；清晰范围自 $H/2 =
\qty{1.9}{m}$ 至无穷远。预设于 $H$ 处，则照相机无需调焦：
可在场景发生的瞬间拍下画面。
\end{solution}

\begin{solution}{exo:b1:optical-instruments:9}
$8\times30$：出瞳 $30/8 = \qty{3.8}{mm} < \qty{7}{mm}$，全部光线
进入眼睛。$7\times50$：\qty{7.1}{mm}，与夜间瞳孔相当，
并多收集 $(50/30)^2 = 2.8$ 倍的光：适合黄昏使用的一对。
$20\times30$：出瞳 \qty{1.5}{mm}，像暗、视场窄，
手的抖动也被放大二十倍。
\end{solution}

\begin{solution}{exo:b1:optical-instruments:10}
隐形眼镜：$f' = \qty{-25}{cm}$，$V = \qty{-4.0}{\delta}$。框架眼镜：
远点距透镜为 $25 - 1.5 = \qty{23.5}{cm}$，因此 $f' =
\qty{-23.5}{cm}$，$V = \qty{-4.3}{\delta}$——光焦度更大。透镜须
使像成于眼睛的远点，该远点与透镜的距离
取决于透镜所在的位置。
\end{solution}

\begin{solution}{exo:b1:optical-instruments:11}
无焦系统：$F_1' = F_2$，间距 $f_1' + f_2' = \qty{20}{cm}$（短于
同放大率 $G$ 的开普勒式望远镜的 \qty{40}{cm}）。$G = -30/(-10) =
+3$：正立。出瞳 $\overline{O_2P'} = f_2'(f_1' + f_2')/f_1' =
-10 \times 20/30 = \qty{-6.7}{cm}$：虚像，位于镜筒内、眼睛无法到达之处；
位于目镜后的眼睛只能接收恰好到达它的光束——
视场狭窄，这是歌剧
望远镜的典型缺点。
\end{solution}

\begin{solution}{exo:b1:optical-instruments:12}
$(200/7)^2 = 820$；$\Delta m = 2.5\log_{10}820 = 7.3$：极限
星等 $13.3$。通量 $\propto 1/r^2$，因此灯可被看见的
距离增至原来的 $\sqrt{820} = 29$ 倍。
\end{solution}

\begin{solution}{pb:b1:optical-instruments:1}
\textbf{1.} $F_2 = F_1'$：间距 $f_1' + f_2' = \qty{1225}{mm}$
（\qty{25}{mm} 目镜）或 \qty{1206}{mm}（\qty{6}{mm}）。

\textbf{2.} $L_1$：$\overline{O_1A_1} = f_1'$（像位于 $F_1'$）。$L_2$：
$\overline{O_2A_1} = -f_2'$，故 $1/\overline{O_2A'} = 1/f_2' + 1/(-f_2')
= 0$：像位于无穷远。

\textbf{3.} 光束在公共焦平面上会聚于高度 $h =
f_1'\alpha$（过 $O_1$ 的光线）；由此处过 $O_2$ 的光线决定
出射方向，$\alpha' = -h/f_2' = -(f_1'/f_2')\alpha$。

\textbf{4.} $G = -48$ 和 $-200$；月球视角：$25^\circ$ 和 $104^\circ$。

\textbf{5.} 天空中没有上下之分；观测望远镜
需加装一对正像棱镜（或一组透镜中继）。

\textbf{6.} 在 $L_1$ 的焦平面上，直径 $f_1'\theta = 1200 \times
9.1\times10^{-3} = \qty{11}{mm}$。

\textbf{7.} 物 $O_1$ 位于 $\overline{O_2O_1} = -(f_1' + f_2')$：
$1/\overline{O_2P'} = 1/f_2' - 1/(f_1' + f_2') = f_1'/[f_2'(f_1' + f_2')]$；
放大率 $\overline{O_2P'}/\overline{O_2O_1} = -f_2'/f_1' = 1/G$，
直径为 $D/\abs G$。

\textbf{8.} \qty{25}{mm}：$\overline{O_2P'} = 25 \times 1225/1200 =
\qty{25.5}{mm}$，直径 $150/48 = \qty{3.1}{mm}$。\qty{6}{mm}：
像位于后方 \qty{6.0}{mm} 处，直径 \qty{0.75}{mm}。

\textbf{9.} 凡进入物镜的光线均经过出瞳，
置于该处的眼睛可接收整个视场；
再向后则只能截取中央光束，视场收缩为锁孔状。

\textbf{10.} $D/\abs G \leq \qty{7}{mm}$：$\abs G \geq 150/7 \approx 21$。

\textbf{11.} $\theta_{\min} = 1.22 \times 5.5\times10^{-7}/0.150 =
\qty{4.5e-6}{rad} = 0.92''$。

\textbf{12.} $4.5\times10^{-6} \times 3.84\times10^5 = \qty{1.7}{km}$；
裸眼为 $3\times10^{-4} \times 3.84\times10^5 = \qty{115}{km}$。

\textbf{13.} $G_r = 3\times10^{-4}/4.5\times10^{-6} = 67$；\qty{25}{mm}
目镜（$48$）最接近，略偏小；\qty{6}{mm}
目镜（$200$）约为所需的三倍。

\textbf{14.} 当 $G = 200$ 时，衍射弥散角为 $200 \times 4.5
\times10^{-6} = \qty{9e-4}{rad} > \epsilon$，肉眼可见地模糊。月球的
光在视网膜上铺开的面积扩大为 $(200/48)^2 = 17$ 倍，
因而更暗。

\textbf{15.} $1'' = \qty{4.85e-6}{rad}$：$D = 1.22\lambda/\theta =
\qty{140}{mm}$。超过此值后，地面观测的分辨率受大气视宁度限制
（不计自适应光学）；更大的 $D$ 仍可收集更多光，
从而能看到更暗的天体。

\textbf{16.} $(150/7)^2 = 460$。

\textbf{17.} $\Delta m = 2.5\log_{10}460 = 6.7$：星等为 $12.7$。

\textbf{18.} $P = 2\times10^{-3} \times \pi(0.075)^2 = \qty{3.5e-5}{W}$，
是裸眼的 $460$ 倍。视网膜像面积扩大为 $\abs G^2$ 倍，
故表面亮度之比为 $460/G^2 = 0.20$（$G = 48$），
$0.012$（$G = 200$）——等于 $(D/\abs G)^2/d_p^2$，
即出瞳与眼瞳直径之比的平方。

\textbf{19.} 对扩展物体，所收集的光随 $D^2$ 增长，
而像面积为 $G^2 \geq (D/d_p)^2$，表面亮度至多
只能等于裸眼值（出瞳 $=$ 眼瞳）。恒星
不可分辨，其全部光落入同一衍射斑中，因此其
亮度随 $(D/d_p)^2$ 增长。

\textbf{20.} 角为 $\beta$ 的点在焦平面上成像，像高 $f_1'\beta$；
光阑允许通过的条件为 $\abs{f_1'\beta} \leq \phi/2$：
全视场为 $\phi/f_1'$（小角近似）。

\textbf{21.} $20/1200 = \qty{0.0167}{rad} = 0.95^\circ$：整个月球
可落入视场内。$5/1200 = 0.24^\circ$：则不能。

\textbf{22.} 表观视场 $\phi/f_2'$：$20/25 = \qty{0.8}{rad} = 46^\circ$；
$5/6 = 48^\circ$。

\textbf{23.} $0.95^\circ = 3420''$：$3420/15 = \qty{230}{s}$，接近四
分钟；$0.24^\circ = 864''$：\qty{58}{s}。

\textbf{24.} \qty{25}{mm} 目镜（$G = 48$，接近 $G_r$，明亮，能看到
整个月面）用于观察圆面；\qty{6}{mm} 目镜（$G = 200$）用于近距离观察环形山，
较暗，也较模糊。仪器可分辨的环形山直径可小至
\qty{1.7}{km}。

\textbf{25.} 焦距 $f_1'$ 与口径 $D$ 相同，则 $G$、视场、出瞳、
分辨率与收光量（减去任何中央遮挡）均相同。优点是无
色差。实际改动是焦点位于主镜
之前，需加一块小副镜，将光束折向
目镜。
\end{solution}
